\section{The exact Ward change on the multiplier range}
\label{sec:short-ward}

Fix the exact real unforced profile $U$ and its full pressure $p$:
\[
 D_0U+U\cdot\nabla U+\nabla p=0,\qquad
 D_\mu=\partial_\tau-\tfrac12y\cdot\nabla-\nu\Delta
                       +\mu-\tfrac12,\qquad \operatorname{div}U=0.
\]
Use normalized phase measure $d\tau/P$, the prescribed orthogonal
symmetry, and the whole-space free inverses
$R_\mu=D_\mu^{-1}$ and $S_\mu=R_\mu\mathbb P\operatorname{div}$.
Here $R_\mu$ acts on divergence-free body forces. Fix
$K\Subset(0,1/4)$ and all geometric parameters. Put
\[
 M_Uw=U\otimes w+w\otimes U,\qquad
 \mathcal D_\mu^T=I+M_US_\mu.
\]
We use only physical stresses $T=M_Uw$ represented in the actual
pair graph of the final return. Its velocity witness $w\in Y_2$
has the stated symmetry and need not be divergence-free. Since $U$ is
bounded, $T\in Y_2$; its old thin part has finite first transverse
moment $N_1$, and its old repair part belongs to
$L^2_\tau L^{22/19}$. The multiplier, moment and repair-product
bounds establish these facts and are preserved by approximation.

\subsection{The full old source and its outward stream cutoff}

Let $r=|y'|$, $s=R^2/r^2$ and $R=\epsilon^{-1}$. Take smooth cutoffs
\[
 \chi_o=0\ (s\le1.2/\nu),\quad \chi_o=1\ (s\ge1.5/\nu),
 \qquad
 \chi_y=1\ (s\le0.8/\nu),\quad \chi_y=0\ (s\ge1/\nu).
\]
Thus the old source disk $r\le R\sqrt{\nu/1.2}$ is separated
from the young receiver $r\ge R\sqrt\nu$. Set
\begin{equation}\label{ward:displacement}
 W_\mu=-S_\mu(\chi_oT),\qquad
 X_0=W_\mu-W_{\mu+1},\qquad
 X=\mathcal C_yX_0,\quad Y=D_{\mu+1}X,\quad h=[U,X].
\end{equation}
Here $[U,X]=U\cdot\nabla X-X\cdot\nabla U$, and the tensor
$\chi_oT=\chi_oM_Vw+\chi_oM_{r_b}w$ is identical in both
resolvents. In particular the full repair stress enters before
inversion and before taking divergence.

For these radially separated divergence-free fields define directly
\[
 \Psi_Z(r,z)=-\int_r^\infty\rho Z_z(\rho,z)\,d\rho,
 \qquad
 \mathcal C_y Z=\chi_yZ+\frac{\chi_y'}r\Psi_Z e_z.
\]
The phase is suppressed. The separated-kernel estimates prove convergence
of this primitive and its differentiated versions for the thin
$N_1$ and repair $L^2_\tau L^{22/19}$ sources. Their decay makes
the boundary term $rZ_r$ at infinity vanish. Hence
$\partial_r\Psi_Z=rZ_z$ and $\partial_z\Psi_Z=-rZ_r$.
The cutoff is divergence-free, equals $Z$ on its plateau and vanishes
in the old region; define it by zero where $\chi_y=0$.
The same kernels control the stream correction and commutators,
so $X,DX,Y\in Y_2$ and, since $U,DU$ are bounded, $h\in Y_2$.
These are fixed-parameter bounds on the actual multiplier range.

The full free resolvent identity gives
\begin{equation}\label{ward:shifted-resolvent}
 X_0=R_{\mu+1}W_\mu,\qquad D_{\mu+1}X_0=W_\mu.
\end{equation}
Consequently the exact cutoff identity is
\begin{equation}\label{ward:cutoff-source}
 Y=\mathcal C_yW_\mu+[D_{\mu+1},\mathcal C_y]X_0.
\end{equation}
It follows first for smooth graph inputs and then by the moment,
repair-product and separated-kernel bounds. For example direct
integration of the outward primitive gives
$\Psi_{AZ}=A\Psi_Z+\Psi_Z$ for $A=-y\cdot\nabla/2$, and
$\Psi_{\Delta Z}=(\partial_r^2-r^{-1}\partial_r+\partial_z^2)\Psi_Z$;
the same decay removes the boundary terms. Thus this identity does
not require a choice of stream primitive near the axis.
Relative to a scalar-cutoff reference return, both
$(\mathcal C_y-\chi_y)W_\mu$ and
$[D_{\mu+1},\mathcal C_y]X_0$ remain. They are supported on the
transition annulus and are separately estimated below.

\subsection{Derivation with the exact pressure}

We derive the identity before projecting. In physical coordinates let
$u$ solve the unforced Navier--Stokes equation with pressure $p_{\rm phys}$,
let $\mathsf X$ be divergence-free, and set
\[
 \mathsf D=\partial_t-\nu\Delta,\quad k(t)=\mu/t,
 \quad \mathsf Y=(\mathsf D+k)\mathsf X,
 \quad \mathsf h=[u,\mathsf X].
\]
Repeated indices are summed, and $(D\mathsf X)_{ij}=\partial_j\mathsf X_i$.
The product rule yields
\[
 (\mathsf D+k)\mathsf h_i
 =[\mathsf Du,\mathsf X]_i+[u,\mathsf Y]_i-2\nu\mathsf C_i,
 \quad
 \mathsf C_i=\sum_{j,l}
 \bigl[(\partial_ju_l)\partial_{jl}\mathsf X_i
       -(\partial_j\mathsf X_l)\partial_{jl}u_i\bigr].
\]
For $N=u\cdot\nabla u$, expansion and cancellation give
\[
 -[N,\mathsf X]_i+u_j\partial_j\mathsf h_i
             +\mathsf h_j\partial_ju_i
       =u_ju_l\partial_{jl}\mathsf X_i.
\]
The pressure identity is
\[
 -[\nabla p_{\rm phys},\mathsf X]_i
 =\partial_i(\mathsf X_j\partial_jp_{\rm phys})
  -(\partial_j\mathsf X_i+\partial_i\mathsf X_j)
                                      \partial_jp_{\rm phys}.
\]
Using $\mathsf Du=-N-\nabla p_{\rm phys}$ and
$[u,\mathsf Y]=-(u\cdot\nabla\mathsf Y+
\mathsf Y\cdot\nabla u)+2u\cdot\nabla\mathsf Y$, these formulas
prove the desired identity. The global Leray projection removes
precisely $\nabla(\mathsf X\cdot\nabla p_{\rm phys})$.

Now substitute
\[
 u=t^{-1/2}U(y,\tau),\quad
 \mathsf X=\sqrt t\,X(y,\tau),\quad
 p_{\rm phys}=t^{-1}p(y,\tau),\quad
 y=x/\sqrt t,\quad\tau=\log t.
\]
Then $\mathsf Y=t^{-1/2}D_{\mu+1}X$ and
$\mathsf h=t^{-1/2}[U,X]$; every term in the differential identity
has factor $t^{-3/2}$. Consequently, in profile variables,
\begin{equation}\label{ward:physical-identity}
 D_\mu h+\mathbb P\operatorname{div}M_Uh
 =-\mathbb P\operatorname{div}M_UY+\mathbb P\widetilde E_\mu,
\end{equation}
where the full unprojected residual is
\begin{align}
 \widetilde E_\mu
   &=2U\cdot\nabla Y-2\nu C_{\rm diff}
       +(U\otimes U):D^2X-(DX+DX^{\mathsf T})\nabla p,
       \label{ward:residual}\\
 (C_{\rm diff})_i
   &=\sum_{j,l}\bigl[(\partial_jU_l)\partial_{jl}X_i
              -(\partial_jX_l)\partial_{jl}U_i\bigr].\notag
\end{align}
All four terms are retained. In particular the pressure is the exact
global pressure of $U$, including the repair. A spatially constant
displacement has zero residual; spatial derivatives of the returned
displacement and velocity produce the remaining force.

For actual graph inputs this identity holds distributionally by
symmetry-preserving smooth approximation of $w$, followed by the
multiplier, moment, repair-product and separated-kernel bounds.
The energy field $h+S_\mu M_U(Y+h)$ has distributional
$D_\mu$ equal to $\mathbb P\widetilde E_\mu$.
The actual-range body estimates below show
$\widetilde E_\mu\in Y_2$: the repair part is bounded directly,
and the thin part differs from its Gaussian affine body by a
$Y_2$ error. Free weak uniqueness therefore gives
\[
 h+S_\mu M_U(Y+h)=R_\mu\mathbb P\widetilde E_\mu.
\]
Multiplying by $M_U$ and adding $M_UY$ yields precisely the
stress identity used in the final return:
\begin{equation}\label{ward:stress-normal}
 \mathcal D_\mu^T M_U(Y+h)
       =M_UY+M_UR_\mu\mathbb P\widetilde E_\mu.
\end{equation}

\subsection{Support and invertibility on the actual range}

Both $Y$ and $h$ are divergence-free and supported in $s\le1/\nu$:
$D_{\mu+1}$ preserves the divergence-free condition, the bracket of divergence-free
fields is divergence-free, and these operations are local after the
stream cutoff. No projection is applied to $Y+h$. Thus
\begin{equation}\label{ward:support-zero}
 \chi_oM_U(Y+h)=0.
\end{equation}
The change $H_oT=M_U(Y[T]+h[T])$ maps the actual multiplier
range to itself by the multiplier bounds. It depends only on
$\chi_oT$ and has young support, so $H_o^2=0$ and its inverse
change is $I-H_o$. The same argument applies after adding the
later affine correction: it depends only on old data and is
supported in $s\le1.1/\nu$, still disjoint from the old source.
No smallness of these changes is required.

In particular a nonzero actual stress $T$ gives a nonzero
$T'=(I+H)T$. If the transformed equation gives
$\mathcal D_\mu^T T'=0$, then $b=-S_\mu T'$ satisfies
$T'=M_Ub$ and is a nonzero divergence-free energy field. This implication
uses neither injectivity of $M_U$ nor an isomorphism of full ambient
velocity and stress kernels. The final return below supplies the
nonzero $T$ and retains both stream transitions, the full pressure,
and all old repair terms in its exact error decomposition.
